\section{Historical appendix: Scriptural Date and Location}
\label{historical-appendix}
% Generated. Do not edit.
% Deterministic projection of research/chronology.toml.
% schema: 5
% document: liturgy/roman-rite/1962/propers/temporal/61-twenty-first-after-pentecost
% generated-by: tools/tpt proper-chronology annotations
\newcommand{\chronologyannotationclaim}[6]{#6}
\newcommand{\chronologyannotationreach}[3]{}
\newcommand{\chronologyannotationgroup}[3]{#3}
\newcommand{\chronologyannotation}[1]{%
  \ifcsname triptychchronologyannotation@#1\endcsname
    \csname triptychchronologyannotation@#1\endcsname
  \else
    \PackageError{triptych}{No chronology annotation for #1}{}%
  \fi}
% introit: Introit
\expandafter\def\csname triptychchronologyannotation@introit\endcsname{%
  \chronologyannotationgroup{historical-setting}{research-pending}{\textbf{Historical setting} -- Occasion date unresolved.}%
  \space \chronologyannotationgroup{composition}{nonuniform}{\textbf{Composition} -- No single date applies across every cited locus.}%
}
% epistle: Epistle
\expandafter\def\csname triptychchronologyannotation@epistle\endcsname{%
  \chronologyannotationgroup{composition}{disputed}{\textbf{Composition} -- disputed: \chronologyannotationreach{Eph.6.10}{Eph}{true}\chronologyannotationreach{Eph.6.11}{Eph}{true}\chronologyannotationreach{Eph.6.12}{Eph}{true}\chronologyannotationreach{Eph.6.13}{Eph}{true}\chronologyannotationreach{Eph.6.14}{Eph}{true}\chronologyannotationreach{Eph.6.15}{Eph}{true}\chronologyannotationreach{Eph.6.16}{Eph}{true}\chronologyannotationreach{Eph.6.17}{Eph}{true}\chronologyannotationclaim{composition.epistle-to-the-ephesians}{composition}{catholic-traditional-v1}{disputed}{a period between 58 and 63}{A.D. 58--63}; \chronologyannotationreach{Eph.6.10}{Eph}{true}\chronologyannotationreach{Eph.6.11}{Eph}{true}\chronologyannotationreach{Eph.6.12}{Eph}{true}\chronologyannotationreach{Eph.6.13}{Eph}{true}\chronologyannotationreach{Eph.6.14}{Eph}{true}\chronologyannotationreach{Eph.6.15}{Eph}{true}\chronologyannotationreach{Eph.6.16}{Eph}{true}\chronologyannotationreach{Eph.6.17}{Eph}{true}\chronologyannotationclaim{composition.epistle-to-the-ephesians}{composition}{catholic-traditional-v1}{disputed}{(Philemon; Colossians; Ephesians; Philippians), 61}{A.D. 61}.}%
}
% gradual: Gradual
\expandafter\def\csname triptychchronologyannotation@gradual\endcsname{%
  \chronologyannotationgroup{historical-setting}{research-pending}{\textbf{Historical setting} -- Occasion date unresolved.}%
  \space \chronologyannotationgroup{composition}{preferred}{\textbf{Composition}: \chronologyannotationreach{Ps.89.1}{Ps}{true}\chronologyannotationreach{Ps.89.2}{Ps}{true}\chronologyannotationclaim{critical.psalms.latest-composition-boundary}{composition}{catholic-critical-v1}{preferred}{before the Maccabean period, around 165 B.C.; no individual psalm can be dated securely}{Before c. 165 B.C.}}%
}
% alleluia: Alleluia
\expandafter\def\csname triptychchronologyannotation@alleluia\endcsname{%
  \chronologyannotationgroup{retrospective-event}{preferred}{\textbf{Retrospective event}: Preferred: \chronologyannotationreach{Ps.113.1}{Ps.77.12-13 Ps.77.42-53 Ps.104.26-38 Ps.104.43 Ps.105.7-12 Ps.105.21-22 Ps.113.1-5 Ps.135.10-15 Ps.80.6-7 2Esd.9.9-12 Acts.7.36}{false}\chronologyannotationclaim{israel.exodus.the-exodus}{retrospective-event}{catholic-traditional-v1}{preferred}{the four hundred and eightieth year after the children of Israel came out of the land of Egypt}{The Exodus from Egypt, the four hundred and eightieth year after the children of Israel came out of the land of Egypt}; alternative: \chronologyannotationreach{Ps.113.1}{Ps.77.12-13 Ps.77.42-53 Ps.104.26-38 Ps.104.43 Ps.105.7-12 Ps.105.21-22 Ps.113.1-5 Ps.135.10-15 Ps.80.6-7 2Esd.9.9-12 Acts.7.36}{false}\chronologyannotationclaim{israel.exodus.the-exodus}{retrospective-event}{catholic-traditional-v1}{alternate}{the Septuagint gives 440 years}{The Exodus from Egypt, the Septuagint gives 440 years}.}%
  \space \chronologyannotationgroup{composition}{preferred}{\textbf{Composition}: \chronologyannotationreach{Ps.113.1}{Ps}{true}\chronologyannotationclaim{critical.psalms.latest-composition-boundary}{composition}{catholic-critical-v1}{preferred}{before the Maccabean period, around 165 B.C.; no individual psalm can be dated securely}{Before c. 165 B.C.}}%
}
% gospel: Gospel
\expandafter\def\csname triptychchronologyannotation@gospel\endcsname{%
  \chronologyannotationgroup{narrated-event}{research-pending}{\textbf{Event} -- No narrated-event date in the chronology corpus.}%
  \space \chronologyannotationgroup{composition}{disputed}{\textbf{Composition} -- disputed: \chronologyannotationreach{Matt.18.23}{Matt}{true}\chronologyannotationreach{Matt.18.24}{Matt}{true}\chronologyannotationreach{Matt.18.25}{Matt}{true}\chronologyannotationreach{Matt.18.26}{Matt}{true}\chronologyannotationreach{Matt.18.27}{Matt}{true}\chronologyannotationreach{Matt.18.28}{Matt}{true}\chronologyannotationreach{Matt.18.29}{Matt}{true}\chronologyannotationreach{Matt.18.30}{Matt}{true}\chronologyannotationreach{Matt.18.31}{Matt}{true}\chronologyannotationreach{Matt.18.32}{Matt}{true}\chronologyannotationreach{Matt.18.33}{Matt}{true}\chronologyannotationreach{Matt.18.34}{Matt}{true}\chronologyannotationreach{Matt.18.35}{Matt}{true}\chronologyannotationclaim{composition.gospel-of-matthew}{composition}{catholic-traditional-v1}{disputed}{about A.D. 38-45}{c. A.D. 38--45}; \chronologyannotationreach{Matt.18.23}{Matt}{true}\chronologyannotationreach{Matt.18.24}{Matt}{true}\chronologyannotationreach{Matt.18.25}{Matt}{true}\chronologyannotationreach{Matt.18.26}{Matt}{true}\chronologyannotationreach{Matt.18.27}{Matt}{true}\chronologyannotationreach{Matt.18.28}{Matt}{true}\chronologyannotationreach{Matt.18.29}{Matt}{true}\chronologyannotationreach{Matt.18.30}{Matt}{true}\chronologyannotationreach{Matt.18.31}{Matt}{true}\chronologyannotationreach{Matt.18.32}{Matt}{true}\chronologyannotationreach{Matt.18.33}{Matt}{true}\chronologyannotationreach{Matt.18.34}{Matt}{true}\chronologyannotationreach{Matt.18.35}{Matt}{true}\chronologyannotationclaim{composition.gospel-of-matthew}{composition}{catholic-traditional-v1}{disputed}{about the year 40-42}{c. A.D. 40--42}; \chronologyannotationreach{Matt.18.23}{Matt}{true}\chronologyannotationreach{Matt.18.24}{Matt}{true}\chronologyannotationreach{Matt.18.25}{Matt}{true}\chronologyannotationreach{Matt.18.26}{Matt}{true}\chronologyannotationreach{Matt.18.27}{Matt}{true}\chronologyannotationreach{Matt.18.28}{Matt}{true}\chronologyannotationreach{Matt.18.29}{Matt}{true}\chronologyannotationreach{Matt.18.30}{Matt}{true}\chronologyannotationreach{Matt.18.31}{Matt}{true}\chronologyannotationreach{Matt.18.32}{Matt}{true}\chronologyannotationreach{Matt.18.33}{Matt}{true}\chronologyannotationreach{Matt.18.34}{Matt}{true}\chronologyannotationreach{Matt.18.35}{Matt}{true}\chronologyannotationclaim{composition.gospel-of-matthew}{composition}{catholic-traditional-v1}{disputed}{the years 40-45}{A.D. 40--45}; \chronologyannotationreach{Matt.18.23}{Matt}{true}\chronologyannotationreach{Matt.18.24}{Matt}{true}\chronologyannotationreach{Matt.18.25}{Matt}{true}\chronologyannotationreach{Matt.18.26}{Matt}{true}\chronologyannotationreach{Matt.18.27}{Matt}{true}\chronologyannotationreach{Matt.18.28}{Matt}{true}\chronologyannotationreach{Matt.18.29}{Matt}{true}\chronologyannotationreach{Matt.18.30}{Matt}{true}\chronologyannotationreach{Matt.18.31}{Matt}{true}\chronologyannotationreach{Matt.18.32}{Matt}{true}\chronologyannotationreach{Matt.18.33}{Matt}{true}\chronologyannotationreach{Matt.18.34}{Matt}{true}\chronologyannotationreach{Matt.18.35}{Matt}{true}\chronologyannotationclaim{composition.gospel-of-matthew}{composition}{catholic-traditional-v1}{disputed}{about the year 60-68}{c. A.D. 60--68}; \chronologyannotationreach{Matt.18.23}{Matt}{true}\chronologyannotationreach{Matt.18.24}{Matt}{true}\chronologyannotationreach{Matt.18.25}{Matt}{true}\chronologyannotationreach{Matt.18.26}{Matt}{true}\chronologyannotationreach{Matt.18.27}{Matt}{true}\chronologyannotationreach{Matt.18.28}{Matt}{true}\chronologyannotationreach{Matt.18.29}{Matt}{true}\chronologyannotationreach{Matt.18.30}{Matt}{true}\chronologyannotationreach{Matt.18.31}{Matt}{true}\chronologyannotationreach{Matt.18.32}{Matt}{true}\chronologyannotationreach{Matt.18.33}{Matt}{true}\chronologyannotationreach{Matt.18.34}{Matt}{true}\chronologyannotationreach{Matt.18.35}{Matt}{true}\chronologyannotationclaim{composition.gospel-of-matthew}{composition}{catholic-traditional-v1}{disputed}{about the years 64-67}{c. A.D. 64--67}; \chronologyannotationreach{Matt.18.23}{Matt}{true}\chronologyannotationreach{Matt.18.24}{Matt}{true}\chronologyannotationreach{Matt.18.25}{Matt}{true}\chronologyannotationreach{Matt.18.26}{Matt}{true}\chronologyannotationreach{Matt.18.27}{Matt}{true}\chronologyannotationreach{Matt.18.28}{Matt}{true}\chronologyannotationreach{Matt.18.29}{Matt}{true}\chronologyannotationreach{Matt.18.30}{Matt}{true}\chronologyannotationreach{Matt.18.31}{Matt}{true}\chronologyannotationreach{Matt.18.32}{Matt}{true}\chronologyannotationreach{Matt.18.33}{Matt}{true}\chronologyannotationreach{Matt.18.34}{Matt}{true}\chronologyannotationreach{Matt.18.35}{Matt}{true}\chronologyannotationclaim{composition.gospel-of-matthew}{composition}{catholic-traditional-v1}{disputed}{about the year 50}{c. A.D. 50}.}%
}
% offertory: Offertory
\expandafter\def\csname triptychchronologyannotation@offertory\endcsname{%
  \chronologyannotationgroup{composition}{preferred}{\textbf{Composition}: \chronologyannotationreach{Job.1.1}{Job}{true}\chronologyannotationreach{Job.1.12}{Job}{true}\chronologyannotationreach{Job.2.6}{Job}{true}\chronologyannotationreach{Job.2.7}{Job}{true}\chronologyannotationclaim{critical.book-of-job}{composition}{catholic-critical-v1}{preferred}{probably sometime between the seventh and fifth centuries B.C.}{Probably sometime between the seventh and fifth centuries B.C.}}%
}
% communion: Communion
\expandafter\def\csname triptychchronologyannotation@communion\endcsname{%
  \chronologyannotationgroup{traditional-attribution}{disputed}{\textbf{Traditional attribution} -- disputed: \chronologyannotationreach{Ps.118.81}{Ps.118}{true}\chronologyannotationreach{Ps.118.84}{Ps.118}{true}\chronologyannotationreach{Ps.118.86}{Ps.118}{true}\chronologyannotationclaim{israel.monarchy.david-traditional-era}{traditional-attribution}{catholic-traditional-v1}{disputed}{reigned from 1055 to 1015 B.C.}{David (reign in the usual chronology), B.C. 1055--1015}.}%
  \space \chronologyannotationgroup{historical-setting}{research-pending}{\textbf{Historical setting} -- Occasion date unresolved.}%
  \space \chronologyannotationgroup{composition}{preferred}{\textbf{Composition}: \chronologyannotationreach{Ps.118.81}{Ps}{true}\chronologyannotationreach{Ps.118.84}{Ps}{true}\chronologyannotationreach{Ps.118.86}{Ps}{true}\chronologyannotationclaim{critical.psalms.latest-composition-boundary}{composition}{catholic-critical-v1}{preferred}{before the Maccabean period, around 165 B.C.; no individual psalm can be dated securely}{Before c. 165 B.C.}}%
}
% presentation: orientation-v1
\newcommand{\chronologyorientationrow}[2]{%
  \ifcsname triptychchronologyorientation@#1\endcsname
    \csname triptychchronologyorientation@#1\endcsname{#2}%
  \else
    \PackageError{triptych}{No chronology orientation row for #1}{}%
  \fi}
\expandafter\def\csname triptychchronologyorientation@row-18bc7e647866f638\endcsname#1{%
Introit: Esth 13:9--11 & Composition: The Book of Esther. #1 & \chronologyannotationreach{Esth.13.10}{Esth}{true}\chronologyannotationreach{Esth.13.11}{Esth}{true}\chronologyannotationreach{Esth.13.9}{Esth}{true}\chronologyannotationclaim{composition.book-of-esther}{composition}{catholic-traditional-v1}{preferred}{at the end of the reign of Xerxes I (485-465 B.C.) or during the reign of his son Artaxerxes I (465-425 B.C.)}{B.C. 465--425 (Conjectural early dating; later date possible)}\\[.25em]%
}
\expandafter\def\csname triptychchronologyorientation@row-d7ec5d94ad1d0e1d\endcsname#1{%
Offertory: Job 1:1, Job 1:12, Job 2:6-7 & Composition: The Book of Job. #1 & \chronologyannotationreach{Job.1.1}{Job}{true}\chronologyannotationreach{Job.1.12}{Job}{true}\chronologyannotationreach{Job.2.6}{Job}{true}\chronologyannotationreach{Job.2.7}{Job}{true}\chronologyannotationclaim{critical.book-of-job}{composition}{catholic-critical-v1}{preferred}{probably sometime between the seventh and fifth centuries B.C.}{Probably sometime between the seventh and fifth centuries B.C.}\\[.25em]%
}
\expandafter\def\csname triptychchronologyorientation@row-b12891d28424f424\endcsname#1{%
Introit: Ps 118:1; Gradual: Psalm 89:1-2; Communion: Psalm 118:81, 84, 86 & Historical setting. #1 & Occasion date unresolved\\[.25em]%
}
\expandafter\def\csname triptychchronologyorientation@row-483ed5cea1be714a\endcsname#1{%
Introit: Ps 118:1; Gradual: Psalm 89:1-2; Alleluia: Psalm 113:1; Communion: Psalm 118:81, 84, 86 & Composition: Psalms. #1 & \chronologyannotationreach{Ps.113.1}{Ps}{true}\chronologyannotationreach{Ps.118.1}{Ps}{true}\chronologyannotationreach{Ps.118.81}{Ps}{true}\chronologyannotationreach{Ps.118.84}{Ps}{true}\chronologyannotationreach{Ps.118.86}{Ps}{true}\chronologyannotationreach{Ps.89.1}{Ps}{true}\chronologyannotationreach{Ps.89.2}{Ps}{true}\chronologyannotationclaim{critical.psalms.latest-composition-boundary}{composition}{catholic-critical-v1}{preferred}{before the Maccabean period, around 165 B.C.; no individual psalm can be dated securely}{Before c. 165 B.C.}\\[.25em]%
}
\expandafter\def\csname triptychchronologyorientation@row-602ada67f8f3bb7d\endcsname#1{%
Alleluia: Psalm 113:1 & Retrospective event: The Exodus from Egypt. #1 & Preferred: \chronologyannotationreach{Ps.113.1}{Ps.77.12-13 Ps.77.42-53 Ps.104.26-38 Ps.104.43 Ps.105.7-12 Ps.105.21-22 Ps.113.1-5 Ps.135.10-15 Ps.80.6-7 2Esd.9.9-12 Acts.7.36}{false}\chronologyannotationclaim{israel.exodus.the-exodus}{retrospective-event}{catholic-traditional-v1}{preferred}{the four hundred and eightieth year after the children of Israel came out of the land of Egypt}{Temple begun in the 480th year after the Exodus (Vulgate)}; alternative: \chronologyannotationreach{Ps.113.1}{Ps.77.12-13 Ps.77.42-53 Ps.104.26-38 Ps.104.43 Ps.105.7-12 Ps.105.21-22 Ps.113.1-5 Ps.135.10-15 Ps.80.6-7 2Esd.9.9-12 Acts.7.36}{false}\chronologyannotationclaim{israel.exodus.the-exodus}{retrospective-event}{catholic-traditional-v1}{alternate}{the Septuagint gives 440 years}{Temple begun in the 440th year after the Exodus (Septuagint)}\\[.25em]%
}
\expandafter\def\csname triptychchronologyorientation@row-b5b7957d62c25282\endcsname#1{%
Introit: Ps 118:1; Communion: Psalm 118:81, 84, 86 & Traditional attribution: David's reign. #1 & Disputed: \chronologyannotationreach{Ps.118.1}{Ps.118}{true}\chronologyannotationreach{Ps.118.81}{Ps.118}{true}\chronologyannotationreach{Ps.118.84}{Ps.118}{true}\chronologyannotationreach{Ps.118.86}{Ps.118}{true}\chronologyannotationclaim{israel.monarchy.david-traditional-era}{traditional-attribution}{catholic-traditional-v1}{disputed}{reigned from 1055 to 1015 B.C.}{B.C. 1055--1015}\\[.25em]%
}
\expandafter\def\csname triptychchronologyorientation@row-88321469d09df574\endcsname#1{%
Gospel: Matthew 18:23-35 & Event. #1 & Date unresolved\\[.25em]%
}
\expandafter\def\csname triptychchronologyorientation@row-02ce0f03aad7eafe\endcsname#1{%
Gospel: Matthew 18:23-35 & Composition: The Gospel of St Matthew. #1 & Disputed: \chronologyannotationreach{Matt.18.23}{Matt}{true}\chronologyannotationreach{Matt.18.24}{Matt}{true}\chronologyannotationreach{Matt.18.25}{Matt}{true}\chronologyannotationreach{Matt.18.26}{Matt}{true}\chronologyannotationreach{Matt.18.27}{Matt}{true}\chronologyannotationreach{Matt.18.28}{Matt}{true}\chronologyannotationreach{Matt.18.29}{Matt}{true}\chronologyannotationreach{Matt.18.30}{Matt}{true}\chronologyannotationreach{Matt.18.31}{Matt}{true}\chronologyannotationreach{Matt.18.32}{Matt}{true}\chronologyannotationreach{Matt.18.33}{Matt}{true}\chronologyannotationreach{Matt.18.34}{Matt}{true}\chronologyannotationreach{Matt.18.35}{Matt}{true}\chronologyannotationclaim{composition.gospel-of-matthew}{composition}{catholic-traditional-v1}{disputed}{about A.D. 38-45}{c. A.D. 38--45}; \chronologyannotationreach{Matt.18.23}{Matt}{true}\chronologyannotationreach{Matt.18.24}{Matt}{true}\chronologyannotationreach{Matt.18.25}{Matt}{true}\chronologyannotationreach{Matt.18.26}{Matt}{true}\chronologyannotationreach{Matt.18.27}{Matt}{true}\chronologyannotationreach{Matt.18.28}{Matt}{true}\chronologyannotationreach{Matt.18.29}{Matt}{true}\chronologyannotationreach{Matt.18.30}{Matt}{true}\chronologyannotationreach{Matt.18.31}{Matt}{true}\chronologyannotationreach{Matt.18.32}{Matt}{true}\chronologyannotationreach{Matt.18.33}{Matt}{true}\chronologyannotationreach{Matt.18.34}{Matt}{true}\chronologyannotationreach{Matt.18.35}{Matt}{true}\chronologyannotationclaim{composition.gospel-of-matthew}{composition}{catholic-traditional-v1}{disputed}{about the year 40-42}{c. A.D. 40--42}; \chronologyannotationreach{Matt.18.23}{Matt}{true}\chronologyannotationreach{Matt.18.24}{Matt}{true}\chronologyannotationreach{Matt.18.25}{Matt}{true}\chronologyannotationreach{Matt.18.26}{Matt}{true}\chronologyannotationreach{Matt.18.27}{Matt}{true}\chronologyannotationreach{Matt.18.28}{Matt}{true}\chronologyannotationreach{Matt.18.29}{Matt}{true}\chronologyannotationreach{Matt.18.30}{Matt}{true}\chronologyannotationreach{Matt.18.31}{Matt}{true}\chronologyannotationreach{Matt.18.32}{Matt}{true}\chronologyannotationreach{Matt.18.33}{Matt}{true}\chronologyannotationreach{Matt.18.34}{Matt}{true}\chronologyannotationreach{Matt.18.35}{Matt}{true}\chronologyannotationclaim{composition.gospel-of-matthew}{composition}{catholic-traditional-v1}{disputed}{the years 40-45}{A.D. 40--45}; \chronologyannotationreach{Matt.18.23}{Matt}{true}\chronologyannotationreach{Matt.18.24}{Matt}{true}\chronologyannotationreach{Matt.18.25}{Matt}{true}\chronologyannotationreach{Matt.18.26}{Matt}{true}\chronologyannotationreach{Matt.18.27}{Matt}{true}\chronologyannotationreach{Matt.18.28}{Matt}{true}\chronologyannotationreach{Matt.18.29}{Matt}{true}\chronologyannotationreach{Matt.18.30}{Matt}{true}\chronologyannotationreach{Matt.18.31}{Matt}{true}\chronologyannotationreach{Matt.18.32}{Matt}{true}\chronologyannotationreach{Matt.18.33}{Matt}{true}\chronologyannotationreach{Matt.18.34}{Matt}{true}\chronologyannotationreach{Matt.18.35}{Matt}{true}\chronologyannotationclaim{composition.gospel-of-matthew}{composition}{catholic-traditional-v1}{disputed}{about the year 60-68}{c. A.D. 60--68}; \chronologyannotationreach{Matt.18.23}{Matt}{true}\chronologyannotationreach{Matt.18.24}{Matt}{true}\chronologyannotationreach{Matt.18.25}{Matt}{true}\chronologyannotationreach{Matt.18.26}{Matt}{true}\chronologyannotationreach{Matt.18.27}{Matt}{true}\chronologyannotationreach{Matt.18.28}{Matt}{true}\chronologyannotationreach{Matt.18.29}{Matt}{true}\chronologyannotationreach{Matt.18.30}{Matt}{true}\chronologyannotationreach{Matt.18.31}{Matt}{true}\chronologyannotationreach{Matt.18.32}{Matt}{true}\chronologyannotationreach{Matt.18.33}{Matt}{true}\chronologyannotationreach{Matt.18.34}{Matt}{true}\chronologyannotationreach{Matt.18.35}{Matt}{true}\chronologyannotationclaim{composition.gospel-of-matthew}{composition}{catholic-traditional-v1}{disputed}{about the years 64-67}{c. A.D. 64--67}; \chronologyannotationreach{Matt.18.23}{Matt}{true}\chronologyannotationreach{Matt.18.24}{Matt}{true}\chronologyannotationreach{Matt.18.25}{Matt}{true}\chronologyannotationreach{Matt.18.26}{Matt}{true}\chronologyannotationreach{Matt.18.27}{Matt}{true}\chronologyannotationreach{Matt.18.28}{Matt}{true}\chronologyannotationreach{Matt.18.29}{Matt}{true}\chronologyannotationreach{Matt.18.30}{Matt}{true}\chronologyannotationreach{Matt.18.31}{Matt}{true}\chronologyannotationreach{Matt.18.32}{Matt}{true}\chronologyannotationreach{Matt.18.33}{Matt}{true}\chronologyannotationreach{Matt.18.34}{Matt}{true}\chronologyannotationreach{Matt.18.35}{Matt}{true}\chronologyannotationclaim{composition.gospel-of-matthew}{composition}{catholic-traditional-v1}{disputed}{about the year 50}{c. A.D. 50}\\[.25em]%
}
\expandafter\def\csname triptychchronologyorientation@row-d944994c29115eb5\endcsname#1{%
Epistle: Ephesians 6:10-17 & Composition: The Epistle to the Ephesians. #1 & Disputed: \chronologyannotationreach{Eph.6.10}{Eph}{true}\chronologyannotationreach{Eph.6.11}{Eph}{true}\chronologyannotationreach{Eph.6.12}{Eph}{true}\chronologyannotationreach{Eph.6.13}{Eph}{true}\chronologyannotationreach{Eph.6.14}{Eph}{true}\chronologyannotationreach{Eph.6.15}{Eph}{true}\chronologyannotationreach{Eph.6.16}{Eph}{true}\chronologyannotationreach{Eph.6.17}{Eph}{true}\chronologyannotationclaim{composition.epistle-to-the-ephesians}{composition}{catholic-traditional-v1}{disputed}{a period between 58 and 63}{A.D. 58--63}; \chronologyannotationreach{Eph.6.10}{Eph}{true}\chronologyannotationreach{Eph.6.11}{Eph}{true}\chronologyannotationreach{Eph.6.12}{Eph}{true}\chronologyannotationreach{Eph.6.13}{Eph}{true}\chronologyannotationreach{Eph.6.14}{Eph}{true}\chronologyannotationreach{Eph.6.15}{Eph}{true}\chronologyannotationreach{Eph.6.16}{Eph}{true}\chronologyannotationreach{Eph.6.17}{Eph}{true}\chronologyannotationclaim{composition.epistle-to-the-ephesians}{composition}{catholic-traditional-v1}{disputed}{(Philemon; Colossians; Ephesians; Philippians), 61}{A.D. 61}\\[.25em]%
}

The date of a narrated event, the composition of a book, and a traditional
attribution are different claims. The following orientation preserves those
distinctions. A book-wide span does not date each verse separately. An
unresolved occasion is left unresolved; the prayers of the Mass are not
assigned biblical dates merely because they echo biblical language. The
rows reproduce the reviewed chronology, with fuller source qualifications
immediately below. They are an orientation to the sources used here, not a
survey of every proposed modern chronology.

\begin{chronologyorientationtable}
\chronologyorientationrow{row-18bc7e647866f638}{Persian court narrative at Susa; the book's place and date of composition remain conjectural.}
\chronologyorientationrow{row-d7ec5d94ad1d0e1d}{The story names Hus (Uz), not the book's place of composition.}
\chronologyorientationrow{row-b12891d28424f424}{No secure individual occasion or place is established for these verses.}
\chronologyorientationrow{row-483ed5cea1be714a}{Collection-wide outer boundary, not one date or place for all four liturgical uses.}
\chronologyorientationrow{row-602ada67f8f3bb7d}{Departure from Egypt remembered in song. The interval ends at the beginning of Solomon's Temple, not at the Psalm's composition.}
\chronologyorientationrow{row-b5b7957d62c25282}{Traditional royal era, conditional upon Davidic attribution; not a dated occasion for Psalm 118.}
\chronologyorientationrow{row-88321469d09df574}{Galilean discourse after the Capharnaum scene and before the departure for Judea; precise teaching venue unspecified.}
\chronologyorientationrow{row-02ce0f03aad7eafe}{Competing historical estimates, some concerning a posited Hebrew or Aramaic original. They do not date the parable's delivery.}
\chronologyorientationrow{row-d944994c29115eb5}{Pauline captivity hypotheses: Rome favored, Caesarea an alternative; recipients in Asia Minor, with a circular-letter theory.}
\end{chronologyorientationtable}

Sources and fuller qualifications: Esther, p.~\pageref{hist-esther};
Job, p.~\pageref{hist-job}; Psalms, p.~\pageref{hist-psalms};
Matthew, p.~\pageref{hist-matthew}; Ephesians, p.~\pageref{hist-ephesians}.

\subsection{Esther: a threatened people in the Persian court}
\label{hist-esther}
The Introit takes its principal words from Mardochai's prayer, not from a
speech of the Persian king. Esther 13:8 identifies the speaker; verses
15--18 identify the threatened people and their appeal to the God who
redeemed them from Egypt. Susa belongs to the narrative setting of the
court (Esth.~1:2;2:5). It is not an established writing desk for the final
book. The historical \work{Catholic Encyclopedia} article ``Esther''
places its conjectural early composition near the end of Xerxes' reign
or in the reign of Artaxerxes, while explicitly allowing later composition
from older materials. The table's compact span is therefore not a secure
date of Mardochai's prayer, and still less a secure date of the chant's
selection. Nor does a hypothesis about the whole book independently date
the Greek prayer material preserved in the Vulgate additions.\footnote{
\work{Catholic Encyclopedia} V (1909), ``Esther,'' authorship and date
discussion, \href{https://www.newadvent.org/cathen/05549a.htm}{retained New Advent article};
Esth.~1:2;2:5;13:8--18, Douay--Rheims/Challoner.}

\subsection{Job: the land of the story and the date of the book}
\label{hist-job}
Hus is the Douay--Rheims form of Uz. Job 1:1 names it but supplies no
chronological synchronism. The reviewed USCCB introduction describes a
probable range for the book's composition spanning the seventh through
fifth centuries before Christ. That is a literary dating judgment, not a
date of Job's losses, nor a settled identification of the place where the
author worked. The Offertory compresses Job 1--2: it joins the initial
portrait, permitted loss of possessions and children, and the second,
bodily trial. Its order must not be flattened into a single uninterrupted
quotation or treated as evidence that the first permission already
included the body.\footnote{Job~1--2; USCCB,
\href{https://bible.usccb.org/bible/job/0}{``Job: Introduction,''}
reviewed source underlying the canonical composition claim. This modern
introduction is cited and paraphrased, not used as study-text English.}

\subsection{The Psalms: remembered events and undated prayers}
\label{hist-psalms}
Psalm 89 bears a Mosaic heading in the biblical witness; Augustine
explicitly cautions that this need not mean Moses wrote it. The Gradual
sings the refuge and eternity clauses without that heading. Psalm 118
does not give these selected verses an independently fixed occasion.
Haydock generally attributes Psalm 118 to David, while reporting three
conjectures: David writing to instruct the young Solomon, David writing
in his youth under Saul's persecution, or Daniel consoling the captives.
The displayed reign follows the historical \work{Catholic Encyclopedia}'s
``usual chronology'' of David. Conditional on Davidic attribution, it dates
neither the Psalm's occasion nor the utterance of verse 81.\footnote{
Augustine, Ps.~89 \S2 (NPNF Ps.~90); Haydock Bible (1859), p.~781,
Ps.~118:1; \work{Catholic Encyclopedia} IV (1908),
\href{https://www.newadvent.org/cathen/04642b.htm}{``David,''} reign chronology.}

The collection-wide outer boundary comes from the reviewed USCCB
Psalms introduction. It does not establish a common composition date for
Psalms 89, 113, and 118, or resolve their individual occasions. The
Alleluia's Exodus is different again: an event remembered by a later
song. The historical ``Biblical Chronology'' article gives the Vulgate
and Septuagint intervals between Exodus and the beginning of the Temple,
discusses substantial chronological difficulties, and proposes its own
reconstruction. The table preserves the two transmitted relative
intervals used in the reviewed chronology, with their actual endpoint.
It does not convert either into an absolute Exodus year or claim that
the Psalm was composed when the Temple was begun.\footnote{USCCB,
\href{https://bible.usccb.org/bible/psalms/0}{``Psalms: Introduction''};
\work{Catholic Encyclopedia} III (1908),
\href{https://www.newadvent.org/cathen/03731a.htm}{``Biblical Chronology,''}
Exodus-to-Temple discussion; 3~Kings 6:1 in Vulgate numbering.}

\subsection{Matthew: the location of the discourse and competing composition dates}
\label{hist-matthew}
Matthew 17:21 places Jesus and the disciples in Galilee; 17:23 names
Capharnaum. Chapter 18 begins ``At that hour,'' and 19:1 narrates the
subsequent departure from Galilee toward Judea beyond Jordan. These
boundaries locate the discourse in the Gospel's narrative movement.
They do not identify a particular room for the parable or assign it a
day or year. The parable itself depicts a king and servants without
supplying their geographical coordinates.

The historical ``St. Matthew'' article reports several competing
composition estimates, including traditions about the apostles'
departure from Palestine and an uncertain interpretation of Irenaeus.
It discusses whether the first written form was Hebrew or Aramaic and
how that relates to the extant Greek Gospel. The Matthew article also reports the early estimate favored by many
Catholic critics of its period; ``New Testament'' supplies the estimate
around the middle of the first century. The table deliberately
retains the distinct estimates rather than averaging them or treating
the earliest as the date of the extant Greek form. They are bounded
historical proposals represented in the reviewed corpus, not a claimed
present consensus. Their disagreement has no power to relocate the
parable within Matthew's own sequence.\footnote{Matt.~17:21--19:1;
\work{Catholic Encyclopedia} X (1911),
\href{https://www.newadvent.org/cathen/10057a.htm}{``St. Matthew,''}
date, language, and composition sections; XIV (1912),
\href{https://www.newadvent.org/cathen/14530a.htm}{``New Testament,''}
origin and composition discussion.}

\subsection{Ephesians: an imprisoned writer and the readers of the letter}
\label{hist-ephesians}
Ephesians 6:20 describes the speaker as an ambassador in a chain.
The historical ``Epistle to the Ephesians'' article reasons within
Pauline authorship and captivity, distinguishing possible imprisonment
at Caesarea from the greater freedom of action associated with Rome.
Its broader interval and the more precise year in ``St. Paul'' are
different historical estimates, not two episodes of composition. The
recipient question is also distinct from the writing-place question:
Ephesus and neighboring communities or a wider circular readership in
Asia Minor are alternatives discussed by the article. The passage does
not say that the armor exhortation was spoken on a battlefield at
Ephesus. Its martial images address Christian endurance and prayer.\footnote{
Eph.~6:18--24; \work{Catholic Encyclopedia} V (1909),
\href{https://www.newadvent.org/cathen/05485a.htm}{``Epistle to the Ephesians,''}
readers and composition discussion; XI (1911),
\href{https://www.newadvent.org/cathen/11567b.htm}{``St. Paul,''}
chronological table and captivity epistles.}
